\documentclass[11pt,a4paper]{article}
\usepackage[margin=1in]{geometry}
\usepackage{times}
\usepackage{amsmath,amssymb}
\usepackage{booktabs}
\usepackage{multirow}
\usepackage{graphicx}
\usepackage{xcolor}
\usepackage{enumitem}
\usepackage[colorlinks=true,citecolor=blue,urlcolor=blue]{hyperref}
\usepackage{natbib}

\newcommand{\Kstar}{K^{*}}
\newcommand{\AUC}{\mathrm{AUC}}

\title{Measuring the Generational Discrimination Boundary\\in Recursively Contaminated Text}

\author{Anonymous}

\begin{document}
\maketitle

% ==============================================================================
\begin{abstract}
Web-scale training corpora are increasingly contaminated by recursively generated synthetic text, yet current detection methods treat synthetic content as a binary category, ignoring the nonlinear quality degradation across generational depth.
We formalize recursive generation as an iterated lossy channel and derive from the Data Processing Inequality (DPI) that a finite generational discrimination boundary $\Kstar$ must exist---the maximum depth at which adjacent generations remain statistically distinguishable.
Through controlled experiments spanning three model families, two generation modes, and cross-model transfer evaluation, we establish $\Kstar=2$ under standard nucleus sampling at the 1.5B parameter scale: human text is readily separable from first-generation synthetic text (pairwise AUC $>0.97$), first- from second-generation text is distinguishable (AUC $\approx 0.61\text{--}0.70$), but second- from third-generation text is not (AUC $<0.56$).
Cross-model transfer experiments show 91\% accuracy retention between model families, indicating that depth signal is predominantly text-intrinsic.
We further demonstrate that $\Kstar$-aware graduated filtering provides modest but consistent improvement over binary filtering for downstream training (up to $+1.1$ percentage points on MMLU at 50\% contamination), though general-purpose data selection methods achieve comparable results, suggesting the primary value of depth measurement lies in scientific characterization.
\end{abstract}

% ==============================================================================
\section{Introduction}
\label{sec:intro}

The proliferation of large language models (LLMs) has led to an accelerating feedback loop: models generate text that enters web corpora, which in turn train future models.
This \emph{recursive contamination} degrades training data quality in ways that binary human-vs-machine detection cannot capture---a first-generation synthetic document may be nearly indistinguishable from human writing in quality, while a fourth-generation document exhibits measurable distributional collapse \cite{shumailov2024collapse}.

Existing approaches address fragments of this problem.
Binary detectors such as DivEye and Binoculars classify text as real or synthetic but do not estimate generational depth.
Model collapse literature describes macroscopic degradation phenomena but does not provide detection tools.
Data selection methods (DSIR, D4, DoReMi) perform distribution matching without explicit depth awareness.
Drayson et al.\ (2025) apply binary detection confidence for importance resampling to mitigate model collapse, representing the closest prior work---yet they too make binary judgments.

We argue that the core missing quantity is the \textbf{generational discrimination boundary} $\Kstar$: the maximum depth $k$ at which generation $k$ and generation $k+1$ remain statistically distinguishable.
From an information-theoretic perspective, recursive generation constitutes an iterated lossy channel; the Data Processing Inequality guarantees that a finite $\Kstar$ must exist.
We provide the first systematic empirical measurement of $\Kstar$.

Our contributions are:
\begin{enumerate}[nosep,leftmargin=1.5em]
    \item \textbf{Characterization of $\Kstar$.} We establish $\Kstar=2$ under standard nucleus sampling at the 1.5B scale across four model configurations, with bootstrap confidence intervals and permutation tests. We show that $\Kstar$ is stable from 1.5B to 7B ($\Delta\mathrm{AUC}_{1v2}=+1.7$pp, $p=0.18$) but sensitive to model variant (base models yield borderline $\Kstar$ at the $\theta=0.60$ threshold). We validate with retrain-chain experiments confirming ecological validity.
    \item \textbf{Cross-model transfer and degradation ordering.} Pythia$\leftrightarrow$OLMo transfer retains 91\% accuracy, demonstrating that depth signal is 70--85\% text-intrinsic. Feature degradation orderings are consistent across models (Spearman $\rho \geq 0.82$).
    \item \textbf{Graduated filtering application.} $\Kstar$-aware graduated filtering yields modest improvement over binary filtering ($+0.6$pp at 30\% contamination, $p=0.08$; $+1.1$pp at 50\%, $p=0.038$), though DSIR achieves comparable performance without explicit depth estimation.
\end{enumerate}

% ==============================================================================
\section{Related Work}
\label{sec:related}

\paragraph{Model collapse and recursive generation.}
Shumailov et al.\ (2024) demonstrate that iteratively training on model-generated data leads to distributional collapse.
Alemohammad et al.\ (2024) and Dohmatob et al.\ (2024) provide theoretical analysis of this phenomenon.
These works characterize macroscopic effects but do not operationalize detection of generational depth in individual documents.

\paragraph{Machine-generated text detection.}
Binary detectors---watermarking, statistical tests (DetectGPT, DNA-GPT), and trained classifiers (DivEye, Binoculars)---distinguish human from machine text but conflate all generational depths into a single ``synthetic'' class.
Drayson et al.\ (2025) use binary detection confidence for importance resampling to prevent model collapse, the closest work to ours; however, we show that binary confidence is a poor depth proxy (Spearman $\rho \leq 0.22$; Section~\ref{sec:binary-proxy}).

\paragraph{Data selection and curation.}
DSIR performs importance resampling to match a target distribution; D4 and DoReMi optimize domain weights.
These methods implicitly down-weight low-quality content but do not explicitly estimate recursive depth.

% ==============================================================================
\section{Method}
\label{sec:method}

\subsection{Recursive Generation as Iterated Lossy Channel}
\label{sec:theory}

Let $P_0$ denote the human text distribution and $P_k = \mathcal{G}(P_{k-1})$ the distribution after $k$ rounds of generation through a stochastic language model $\mathcal{G}$.
Each application of $\mathcal{G}$ constitutes a lossy channel: the generated distribution $P_k$ cannot contain more information about $P_0$ than $P_{k-1}$ does.
By the Data Processing Inequality, for any divergence measure $D$:
\begin{equation}
    D(P_k \| P_{k+1}) \leq D(P_{k-1} \| P_k) \leq \cdots \leq D(P_0 \| P_1).
    \label{eq:dpi}
\end{equation}

Since $D(P_k \| P_{k+1})$ is monotonically non-increasing and bounded below by zero, there exists a finite $\Kstar_{\mathrm{theory}}$ beyond which adjacent generations are information-theoretically indistinguishable.
Our empirical $\Kstar_{\mathrm{empirical}}$ provides a conservative lower bound on $\Kstar_{\mathrm{theory}}$, as measurement noise raises the effective indistinguishability floor.

\subsection{Depth Estimation via Distributional Features}
\label{sec:features}

We extract a 15-dimensional hand-crafted feature vector from each document, designed to capture distributional shifts induced by recursive generation:

\begin{itemize}[nosep,leftmargin=1.5em]
    \item \textbf{Surprisal features} (6D): mean, std, and first/second differences of per-token surprisal under a reference language model, capturing smoothing of the surprisal landscape across generations.
    \item \textbf{Lexical diversity features} (5D): type-token ratio, hapax ratio, vocabulary richness, frequency entropy, and frequency kurtosis, capturing vocabulary compression.
    \item \textbf{Structural features} (4D): sentence length statistics, discourse markers, and repetition metrics, capturing homogenization of text structure.
\end{itemize}

We train a LightGBM ordinal classifier on these features, predicting depth class $\hat{k} \in \{0, 1, 2, \ldots\}$.

\subsection{Measuring $\Kstar$}
\label{sec:kstar-def}

We operationalize $\Kstar$ as the largest $k$ for which the pairwise AUC between generations $k$ and $k+1$ exceeds a threshold $\theta$:
\begin{equation}
    \Kstar = \max\{k : \AUC(k, k+1) > \theta\}.
    \label{eq:kstar}
\end{equation}

We set $\theta = 0.60$ (moderate discriminability above the 0.50 chance level) and report bootstrap 95\% confidence intervals on each pairwise AUC.
$\Kstar = k$ requires: (i) the lower bound of the 95\% CI for $\AUC(k-1, k)$ exceeds $\theta$, and (ii) the upper bound of the 95\% CI for $\AUC(k, k+1)$ falls below $\theta$.

\subsection{Graduated Filtering}
\label{sec:graduated}

Given a predicted depth $\hat{k}$ and a mixing parameter $\alpha \in [0,1]$, we assign sample weight:
\begin{equation}
    w(\hat{k}) = \begin{cases}
    1 & \text{if } \hat{k} = 0, \\
    (1 - \alpha)^{\hat{k}} & \text{if } \hat{k} \geq 1.
    \end{cases}
    \label{eq:graduated}
\end{equation}

This provides a calibrated weighting that retains high-quality first-generation text at reduced weight while aggressively down-weighting deeper generations, in contrast to binary filtering which discards all synthetic text uniformly.

% ==============================================================================
\section{Characterization Experiments}
\label{sec:characterization}

All experiments use continuation-mode generation as the canonical condition unless otherwise noted, following our analysis that rewrite (instruction) mode introduces mode-specific style artifacts that inflate depth signal but reduce cross-model transferability (Section~\ref{sec:genmode}).

\subsection{$\Kstar$ Measurement}
\label{sec:kstar-measurement}

\paragraph{Setup.}
We measure $\Kstar$ across four model configurations using C4 documents, generating recursive chains to depth 5 with nucleus sampling ($p=0.95$, $T=1.0$).
Table~\ref{tab:kstar-main} reports pairwise AUC with bootstrap 95\% CIs.

\begin{table}[t]
\centering
\small
\caption{Pairwise AUC across model configurations (continuation mode, $\theta=0.60$). All configurations yield $\Kstar=2$: 1v2 CI lower $> 0.60$, 2v3 CI upper $< 0.60$.}
\label{tab:kstar-main}
\begin{tabular}{lcccc}
\toprule
\textbf{Config} & \textbf{0v1} & \textbf{1v2} & \textbf{2v3} & $\Kstar$ \\
\midrule
Qwen-1.5B (rw) & 0.996 & 0.696 & 0.565 & 2 \\
Pythia-1.4B & 0.993 & 0.654 & 0.552 & 2 \\
OLMo-1B & 0.985 & 0.615 & 0.548 & 2 \\
Qwen-1.5B (cont) & 0.992 & 0.611 & 0.548 & 2 \\
\bottomrule
\end{tabular}
\end{table}

Across all four configurations, $\Kstar=2$.
The 0v1 boundary is consistently strong (AUC $> 0.98$), while the 1v2 boundary ranges from $0.611$ to $0.696$, with the continuation-mode Qwen configuration being most marginal (CI: $[0.602, 0.621]$).
All 2v3 AUCs fall below $0.57$, confirming that depth $\geq 2$ text converges to a distributional fixed point.

\paragraph{Cross-scale stability.}
Table~\ref{tab:cross-scale} compares Qwen2.5-1.5B and 7B.
The 7B model yields $\AUC_{1v2} = 0.629$ versus $0.611$ for 1.5B ($\Delta = +1.7$pp, Welch $t$-test $p = 0.18$), a non-significant difference.

\begin{table}[t]
\centering
\small
\caption{Cross-scale comparison: Qwen2.5 1.5B vs 7B (continuation mode, nucleus $p=0.95$). $\Kstar=2$ remains stable; scale does not significantly affect the discrimination boundary.}
\label{tab:cross-scale}
\begin{tabular}{lccccc}
\toprule
\textbf{Scale} & \textbf{0v1} & \textbf{1v2} & \textbf{1v2 CI} & \textbf{2v3} & $p$ \\
\midrule
1.5B & 0.992 & 0.611 & [.602, .621] & 0.548 & --- \\
7B & 0.998 & 0.629 & [.606, .651] & 0.558 & 0.18 \\
\bottomrule
\end{tabular}
\end{table}

The wider CI for 7B ($\pm 0.023$) reflects the reduced sample size (1K vs 5K documents) necessitated by computational constraints.
The degradation ordering between scales is highly consistent (Spearman $\rho = 0.82$), with \texttt{surp\_d1\_std}, \texttt{surp\_std}, and \texttt{freq\_entropy} as the top discriminative features at both scales.

\paragraph{Base vs instruct model.}
Table~\ref{tab:base-instruct} compares Qwen2.5-1.5B base (no instruction tuning) against the instruct variant.
The base model yields $\AUC_{1v2} = 0.599$, with CI lower bound $0.588 < 0.60$---placing $\Kstar$ at the threshold boundary.

\begin{table}[t]
\centering
\small
\caption{Base vs instruct model (Qwen2.5-1.5B, continuation mode). The base model's 1v2 AUC CI lower falls below $\theta=0.60$, yielding borderline $\Kstar$.}
\label{tab:base-instruct}
\begin{tabular}{lcccc}
\toprule
\textbf{Variant} & \textbf{0v1} & \textbf{1v2} & \textbf{1v2 CI} & $\Kstar$ \\
\midrule
Instruct & 0.992 & 0.611 & [.602, .621] & 2 \\
Base & 0.992 & 0.599 & [.588, .609] & 1\textsuperscript{$\dagger$} \\
\bottomrule
\end{tabular}
\vspace{0.3em}
\raggedright\scriptsize\textsuperscript{$\dagger$}At $\theta=0.585$, the base model also achieves $\Kstar=2$; the result is threshold-sensitive.
\end{table}

Feature degradation orderings remain highly consistent between base and instruct variants (Spearman $\rho = 0.87$), with \texttt{surp\_d1\_std} and \texttt{surp\_std} ranking in the top 3 for both.
The key divergence is \texttt{freq\_entropy}, which ranks 1st for the instruct model but 5th for the base model, reflecting instruction tuning's effect on vocabulary concentration.

\subsection{Cross-Model Transfer}
\label{sec:transfer}

Depth features trained on one model's generations and evaluated on another's retain 91\% accuracy (Pythia$\leftrightarrow$OLMo), with 0v1 transfer AUC exceeding 0.96 across all pairs.
This demonstrates that depth signal is predominantly text-intrinsic (70--85\%) rather than an artifact of model self-familiarity.
Transfer performance for 1v2+ pairs depends on model family similarity, consistent with the weak signal above $\Kstar$.

\subsection{Degradation Ordering}
\label{sec:degradation}

We analyze which features degrade most rapidly across generations.
Single-feature AUC rankings are consistent across model configurations (Spearman $\rho \geq 0.82$ for continuation mode).
The top three discriminative features---\texttt{surp\_d1\_std} (first-difference surprisal standard deviation), \texttt{surp\_std} (raw surprisal standard deviation), and \texttt{surp\_d2\_std} or \texttt{freq\_entropy}---reflect the smoothing of the surprisal landscape and compression of vocabulary diversity that characterize iterative generation.

The retrain-chain experiment (Section~\ref{sec:retrain}) confirms that this ordering transfers to the ecologically valid setting ($\rho = 0.78$, with 2 of 3 top features shared), with \texttt{freq\_kurtosis} gaining importance in the retrain scenario due to LoRA fine-tuning's systematic compression of frequency distribution tails.

\subsection{Generation Mode Analysis}
\label{sec:genmode}

Rewrite (instruction) mode yields substantially higher signal (AUC$_{1v2} = 0.696$ vs $0.611$ for continuation, $\Delta = +8.4$pp) but lower cross-model transferability (61--63\% retention vs 80--92\%).
The asymmetry in feature transfer---continuation$\to$rewrite 89\% vs rewrite$\to$continuation 82\%---indicates that continuation features form a universal subset augmented by mode-specific style artifacts in rewrite mode.
We therefore adopt continuation mode as the canonical condition for all cross-model claims.

% ==============================================================================
\section{Application: Graduated Filtering}
\label{sec:application}

\subsection{Filtering Strategy Comparison}
\label{sec:filtering}

We evaluate six filtering strategies under continued pretraining of Qwen2.5-1.5B with 30\% synthetic text injection, measuring downstream performance on MMLU and HellaSwag across 5 random seeds.
Table~\ref{tab:filtering} reports results.

\begin{table}[t]
\centering
\small
\caption{Filtering strategies at 30\% synthetic injection (5 seeds, mean $\pm$ std). Graduated filtering yields modest improvement over binary; DSIR achieves comparable performance without depth estimation.}
\label{tab:filtering}
\begin{tabular}{lcc}
\toprule
\textbf{Strategy} & \textbf{MMLU} & \textbf{HellaSwag} \\
\midrule
Clean baseline & 46.5 $\pm$ 0.3 & 62.0 $\pm$ 0.4 \\
\midrule
No filter & 44.8 $\pm$ 0.4 & 60.2 $\pm$ 0.5 \\
Surprisal-based & 45.0 $\pm$ 0.4 & 60.4 $\pm$ 0.5 \\
Binary (Drayson) & 45.3 $\pm$ 0.3 & 60.8 $\pm$ 0.4 \\
Binary (ours) & 45.5 $\pm$ 0.3 & 60.9 $\pm$ 0.4 \\
DSIR & 45.9 $\pm$ 0.3 & 61.4 $\pm$ 0.4 \\
Graduated ($\alpha\!=\!0.5$) & \textbf{46.1} $\pm$ 0.3 & \textbf{61.5} $\pm$ 0.4 \\
\bottomrule
\end{tabular}
\end{table}

Graduated filtering recovers most of the clean-data performance ($-0.4$pp on MMLU vs clean, compared to $-1.7$pp for no filtering).
However, the advantage over binary filtering is modest: $+0.6$pp on both MMLU and HellaSwag, with paired $t$-test $p$-values of $0.072$ and $0.085$ respectively---borderline significant.
DSIR performs comparably to graduated filtering ($-0.2$pp MMLU, $-0.1$pp HellaSwag, neither significant), indicating that importance resampling implicitly achieves a similar quality gradient without explicit depth estimation.

\paragraph{Sensitivity to $\alpha$.}
Table~\ref{tab:alpha} shows that graduated filtering is robust to the mixing parameter, with a shallow inverted-U peaking at $\alpha = 0.5$.

\begin{table}[t]
\centering
\small
\caption{Graduated filtering sensitivity to $\alpha$ (30\% injection).}
\label{tab:alpha}
\begin{tabular}{lccccc}
\toprule
$\alpha$ & 0.1 & 0.3 & 0.5 & 0.7 & 0.9 \\
\midrule
MMLU & 45.7 & 45.9 & \textbf{46.1} & 46.0 & 45.8 \\
HellaSwag & 61.0 & 61.3 & \textbf{61.5} & 61.4 & 61.1 \\
\bottomrule
\end{tabular}
\end{table}

\subsection{Dose-Response Analysis}
\label{sec:dose}

The practical value of graduated filtering scales with contamination prevalence.
Table~\ref{tab:dose} shows the graduated$-$binary delta across three injection rates.

\begin{table}[t]
\centering
\small
\caption{Dose-response: graduated vs binary filtering across contamination levels. Benefit becomes significant only at high contamination (50\%).}
\label{tab:dose}
\begin{tabular}{lccccc}
\toprule
\textbf{Dose} & \multicolumn{2}{c}{\textbf{MMLU}} & \multicolumn{2}{c}{\textbf{HellaSwag}} & $p$ \\
\cmidrule(lr){2-3}\cmidrule(lr){4-5}
 & Binary & Grad. & Binary & Grad. & (MMLU) \\
\midrule
10\% & 46.2 & 46.4 & 61.7 & 61.9 & 0.32 \\
30\% & 45.5 & 46.1 & 60.9 & 61.5 & 0.08 \\
50\% & 44.8 & 45.9 & 60.1 & 61.0 & 0.038 \\
\bottomrule
\end{tabular}
\end{table}

At 10\% contamination, the difference is negligible ($+0.2$pp, $p=0.32$).
At 50\% contamination---a worst-case but increasingly plausible scenario as web synthetic content grows---the benefit reaches $+1.1$pp on MMLU ($p=0.038$).

% ==============================================================================
\section{Supporting Analyses}
\label{sec:supporting}

\subsection{Retrain-Chain Ecological Validity}
\label{sec:retrain}

Prompt-chain generation provides controlled conditions but may overestimate inter-generational differences compared to the recursive fine-tuning that characterizes real-world model collapse.
We validate with a retrain-chain experiment: Pythia-1.4B with LoRA (rank=16, $\alpha=32$) iteratively fine-tuned on its own generations (C4, 5K documents, depth 0--4).

\begin{table}[t]
\centering
\small
\caption{Prompt-chain vs retrain-chain pairwise AUC (Pythia-1.4B). Retrain-chain confirms $\Kstar=2$ with attenuated but consistent signal.}
\label{tab:retrain}
\begin{tabular}{lcccc}
\toprule
\textbf{Method} & \textbf{0v1} & \textbf{1v2} & \textbf{2v3} & \textbf{3v4} \\
\midrule
Prompt-chain & 0.993 & 0.654 & 0.552 & --- \\
Retrain-chain & 0.975 & 0.621 & 0.544 & 0.510 \\
$\Delta$ & $-$1.8pp & $-$3.3pp & $-$0.8pp & --- \\
\bottomrule
\end{tabular}
\end{table}

The retrain-chain attenuates all pairwise AUCs ($\Delta\AUC_{1v2} = -3.3$pp) but preserves $\Kstar=2$.
Feature ranking correlation between prompt- and retrain-chains is $\rho = 0.78$, with 2 of 3 top features shared.
The attenuation reflects that LoRA fine-tuning produces more subtle distributional shifts than prompt-based generation, as the model gradually adapts to its own output distribution.

\subsection{LLM Zero-Shot Depth Estimation}
\label{sec:llm-baseline}

We evaluate GPT-4o and Claude 3.5 on zero-shot depth estimation to establish that generational depth is a distributional property not accessible to surface-level analysis.
Table~\ref{tab:llm-zeroshot} reports results.

\begin{table}[t]
\centering
\small
\caption{LLM zero-shot vs our 15D method on depth estimation. LLMs achieve near-chance on depth classification despite reasonable binary detection.}
\label{tab:llm-zeroshot}
\begin{tabular}{lccccc}
\toprule
\textbf{Method} & \textbf{Bin.\ Acc.} & \textbf{3-cl.} & \textbf{6-cl.} & $\rho$ & \textbf{MAE} \\
\midrule
GPT-4o & 0.72 & 0.38 & 0.20 & 0.18 & 1.85 \\
Claude 3.5 & 0.68 & 0.36 & 0.19 & 0.15 & 1.92 \\
\midrule
Random & --- & 0.33 & 0.17 & --- & --- \\
Ours (15D) & --- & 0.63 & 0.42 & 0.55 & --- \\
\bottomrule
\end{tabular}
\end{table}

While LLMs achieve moderate binary detection (68--72\%), their 3-class depth accuracy (36--38\%) is near chance (33.3\%), and Spearman $\rho$ with true depth is $\leq 0.18$.
Our 15D feature method achieves 63\% 3-class accuracy and $\rho = 0.55$, confirming that depth estimation requires specialized distributional measurement.

\subsection{Binary Detector Confidence as Depth Proxy}
\label{sec:binary-proxy}

A natural question is whether binary detector confidence scores can serve as depth proxies, obviating the need for dedicated depth estimation.
Table~\ref{tab:binary-proxy} shows this is not the case.

\begin{table}[t]
\centering
\small
\caption{Binary detector confidence as depth proxy. Strong binary detection does not translate to depth estimation.}
\label{tab:binary-proxy}
\begin{tabular}{lcccc}
\toprule
\textbf{Detector} & \textbf{Bin.\ AUC} & $\rho_{\mathrm{depth}}$ & \textbf{AUC}$_{1v2}$ & \textbf{AUC}$_{2v3}$ \\
\midrule
DivEye & 0.882 & 0.215 & 0.548 & 0.512 \\
Binoculars & 0.907 & 0.192 & 0.534 & 0.508 \\
\midrule
Ours (15D) & --- & 0.553 & 0.654 & 0.552 \\
\bottomrule
\end{tabular}
\end{table}

Binary confidence correlates weakly with depth ($\rho \leq 0.22$) and yields near-chance pairwise AUC for 1v2 and 2v3 pairs.
Notably, Binoculars achieves higher binary AUC (0.907) but \emph{lower} depth correlation ($\rho = 0.192$) than DivEye ($\rho = 0.215$), confirming that better binary detection does not imply better depth estimation---these are fundamentally different tasks.

\subsection{Learned Representation Baseline}
\label{sec:learned}

We compare our hand-crafted 15D features against frozen Qwen2.5-1.5B last-hidden-state embeddings (1536D) with a 2-layer MLP, evaluated across three generator configurations (Table~\ref{tab:learned}).

\begin{table}[t]
\centering
\small
\caption{Hand-crafted 15D features vs frozen LM embeddings + MLP. Learned representations outperform by $\sim$5pp but at the cost of interpretability.}
\label{tab:learned}
\begin{tabular}{llccc}
\toprule
\textbf{Generator} & \textbf{Method} & \textbf{Acc.} & \textbf{mPW-AUC} & \textbf{AUC}$_{1v2}$ \\
\midrule
\multirow{2}{*}{Qwen (rw)} & 15D & 0.423 & 0.621 & 0.696 \\
 & Learned & 0.479 & 0.682 & 0.734 \\
\midrule
\multirow{2}{*}{Pythia} & 15D & 0.386 & 0.606 & 0.654 \\
 & Learned & 0.436 & 0.650 & 0.697 \\
\midrule
\multirow{2}{*}{OLMo} & 15D & 0.382 & 0.602 & 0.615 \\
 & Learned & 0.431 & 0.643 & 0.658 \\
\bottomrule
\end{tabular}
\end{table}

Learned representations outperform hand-crafted features by $\sim$5pp in accuracy and $\sim$4.9pp in mean pairwise AUC, consistently across all configurations.
However, the 15D features capture approximately 85\% of the available discriminative signal, indicating that depth information is predominantly encoded in low-dimensional statistical patterns.
The interpretability-performance tradeoff is central: learned features are stronger but opaque, while hand-crafted features uniquely support the degradation ordering analysis (Section~\ref{sec:degradation}) and cross-model transfer analysis (Section~\ref{sec:transfer}) that constitute the paper's primary characterization contribution.

\subsection{Naturalistic Validation}
\label{sec:naturalistic}

We apply our depth estimator to a 10K-document sample from Common Crawl (2025 snapshot) to assess ecological plausibility.

\begin{table}[t]
\centering
\small
\caption{Predicted depth distribution on Common Crawl 2025 and correlation with independent quality proxies.}
\label{tab:naturalistic}
\begin{tabular}{lclc}
\toprule
\textbf{Pred.\ Depth} & \textbf{\%} & \textbf{Proxy} & $\rho$ \\
\midrule
Gen-0 (human) & 83.2 & AlignScore & $-$0.28 \\
Gen-1 & 11.5 & Vocab richness & $-$0.24 \\
Gen-2+ & 5.3 & Self-BLEU & +0.19 \\
 & & Perplexity (GPT-2) & $-$0.31 \\
\bottomrule
\end{tabular}
\end{table}

The predicted distribution (83/12/5\%) is broadly consistent with current estimates of web synthetic content prevalence (10--20\%).
Correlations with independent quality proxies are weak but directionally consistent: predicted deeper-generation text exhibits lower factual alignment (AlignScore $\rho = -0.28$), lower vocabulary richness ($\rho = -0.24$), and lower perplexity under GPT-2 ($\rho = -0.31$), the latter reflecting proximity to the LM typical set.
We present this as preliminary ecological validation, acknowledging the absence of ground-truth depth labels and the confounding effects of domain, length, and genre variation.

% ==============================================================================
\section{Conclusion}
\label{sec:conclusion}

We introduced the generational discrimination boundary $\Kstar$ as a principled quantity for characterizing recursive synthetic text contamination, grounded in the Data Processing Inequality.
Through systematic measurement across model families, scales, and generation modes, we established $\Kstar=2$ under standard nucleus sampling at the 1.5B scale---meaning that while human and first-generation synthetic text are readily distinguishable, text beyond the second generation converges to a distributional fixed point.
This finding is stable from 1.5B to 7B parameters ($p=0.18$) and across model families (91\% cross-model transfer retention), though base models without instruction tuning show borderline behavior at the threshold.

The practical implication is nuanced: $\Kstar$-aware graduated filtering provides modest improvement over binary filtering, scaling with contamination prevalence ($+1.1$pp at 50\% contamination), but general-purpose data selection achieves comparable results.
The primary value of $\Kstar$ measurement thus lies in scientific characterization---understanding \emph{how} recursive generation degrades text distributions---rather than direct engineering advantage.

\paragraph{Limitations.}
Our measurements are limited to 1.5B--7B scale models; larger models may exhibit different $\Kstar$ behavior.
The retrain-chain validation uses LoRA (rank=16) rather than full fine-tuning.
Naturalistic validation lacks ground-truth depth labels.
The $\Kstar=2$ result is specific to nucleus sampling ($p=0.95$); other decoding strategies (e.g., greedy) may yield different boundaries.

\end{document}
